% Source: Weinberg GR, physical PDF pp. 386--394
%         (printed pp. 365--373).
% Coverage: Section 12.5, The Tetrad Formalism; equations
%           (12.5.1)--(12.5.36).
% Figures/tables/footnotes: one optional-reading footnote; reference marker 1;
%                          no figures or tables.
% Status: source-reviewed and compile-clean.
% Uncertainties: tetrad notation and the source's inverse-tetrad/co-tetrad
%                variation-sign mismatch are documented below; the co-tetrad
%                formulation is checked against the Hilbert stress tensor.
%
% MODERNIZATION CHECK: Weinberg's V^alpha_mu is e^a_mu, its inverse is
% e_a^mu, local-frame indices are Latin, and his starred local components are
% written directly with Latin indices. Partial/covariant derivatives are
% compact, and coordinate displacements are xi.
%
% SOURCE ANOMALY / MODERNIZATION CHOICE: In (12.5.26) the source varies the
% inverse tetrad V_alpha^mu, but (12.5.27) assigns U the sign appropriate to a
% variation of the co-tetrad. Taken literally, those signs conflict with the
% metric variation printed before (12.5.35). Equations (12.5.26)--(12.5.35)
% below consistently use the co-tetrad e^a_mu as the varied field. This keeps
% Weinberg's positive definition (12.5.27), the binding Hilbert definition
% T_mn=-2/sqrt(-g) delta I_M/delta g^mn, and (12.5.36) mutually consistent.
%
% LITERAL-OPERATOR AUDIT: every continuation-line operator in
% (12.5.14), (12.5.17), (12.5.22)--(12.5.25), and
% (12.5.32)--(12.5.35) was checked against physical PDF pp. 389--393.

\subsection{The Tetrad Formalism\optionalreading}\label{sec:12.5}
\begingroup
\renewcommand{\thefootnote}{*}
\begin{NoHyper}
\footnotetext{This section lies somewhat out of the book's main line of
development, and may be omitted in a first reading.}
\end{NoHyper}
\endgroup

Until now, we have followed only one approach in determining the effects of
gravitation on general physical systems. Given the special-relativistic
equations that govern the system in the absence of gravitation, we replace
all Lorentz tensors with objects that behave like tensors (or tensor
densities) under general coordinate transformations. Also, we replace all
ordinary derivatives with covariant derivatives, and replace
\(\eta_{ab}\) everywhere with \(g_{\mu\nu}\). The equations of motion are
then generally covariant. (See Chapter~5.)

This method actually works \emph{only} for objects that behave like tensors
under Lorentz transformation, and not for the spinor fields discussed in
Section~\ref{sec:2.12}. (Mathematically, this is because the tensor
representations of the group \(\mathrm{GL}(4)\) of general linear
\(4\times4\) matrices behave like tensors under the subgroup of Lorentz
transformations, but there are no representations of \(\mathrm{GL}(4)\), or
even ``representations up to a sign,'' which behave like spinors under the
Lorentz subgroup.) How then can we incorporate spinors into general
relativity?

The answer lies in a different approach to the problem of determining the
effects of gravitation on physical systems, an approach that is rather
interesting in its own right, quite apart from the problem of dealing with
spinors.

To start, let us take advantage of the Principle of Equivalence, and at every
point \(X\) erect a set of coordinates \(\xi_X^a\) that are locally inertial
at \(X\). (Of course, it will not be possible to erect any \emph{single}
coordinate system that is locally inertial everywhere, unless the space-time
continuum is ``flat.'') As shown in Sections~\ref{sec:3.2} and \ref{sec:3.3},
the metric in any general noninertial coordinate system is then
\begin{equation}
  g_{\mu\nu}(x)
  =
  e^a{}_\mu(x)e^b{}_\nu(x)\eta_{ab},
  \tag{12.5.1}\label{eq:12.5.1}
\end{equation}
where
\begin{equation}
  e^a{}_\mu(X)
  \equiv
  \left.
  \frac{\partial\xi_X^a(x)}{\partial x^\mu}
  \right|_{x=X}.
  \tag{12.5.2}\label{eq:12.5.2}
\end{equation}
Note that we fix the locally inertial coordinates \(\xi_X^a\) once and for
all at each physical point \(X\), so when we change our general noninertial
coordinates from \(x^\mu\) to \(x'^\mu\), the partial derivatives
\(e^a{}_\mu\) change according to the rule
\begin{equation}
  e^a{}_\mu
  \longrightarrow
  e'^a{}_\mu
  =
  \frac{\partial x^\nu}{\partial x'^\mu}e^a{}_\nu.
  \tag{12.5.3}\label{eq:12.5.3}
\end{equation}
Thus, we are to think of \(e^a{}_\mu\) as forming \emph{four} covariant
vector fields, \emph{not} as a single coordinate tensor with two coordinate
indices. This set of four vectors is known as a \emph{tetrad}, or
\emph{vierbein}.

Now, given any contravariant vector field \(A^\mu(x)\), we can use the tetrad
to refer its components at \(x\) to the coordinate system \(\xi_x^a\)
locally inertial at \(x\):
\begin{equation}
  A^a\equiv e^a{}_\mu A^\mu.
  \tag{12.5.4}\label{eq:12.5.4}
\end{equation}
We are contracting a contravariant vector \(A^\mu\) with four covariant
vectors \(e^a{}_\mu\), so this has the effect of replacing the single
four-vector \(A^\mu\) with the four coordinate scalars \(A^a\). We can do the
same with covariant vector fields, and indeed with general tensor fields:
\begin{equation}
  A_a\equiv e_a{}^\mu A_\mu,
  \qquad
  B^a{}_b\equiv e^a{}_\mu e_b{}^\nu B^\mu{}_\nu,
  \quad\text{etc.}
  \tag{12.5.5}\label{eq:12.5.5}
\end{equation}
Here \(e_a{}^\mu\) is just the tetrad \eqref{eq:12.5.2}, but with the local
index lowered with the Minkowski tensor and the coordinate index raised with
the metric tensor:
\begin{equation}
  e_a{}^\mu
  \equiv
  \eta_{ab}g^{\mu\nu}e^b{}_\nu.
  \tag{12.5.6}\label{eq:12.5.6}
\end{equation}
Note that according to Eq.~\eqref{eq:12.5.1} this is just the inverse of the
tetrad:
\begin{equation}
  \delta^\mu{}_\nu=e_a{}^\mu e^a{}_\nu,
  \tag{12.5.7}\label{eq:12.5.7}
\end{equation}
and hence also
\begin{equation}
  \delta^a{}_b=e^a{}_\mu e_b{}^\mu.
  \tag{12.5.8}\label{eq:12.5.8}
\end{equation}
The coordinate-scalar components of the metric tensor are then simply
\begin{equation}
  g_{ab}
  \equiv
  e_a{}^\mu e_b{}^\nu g_{\mu\nu}
  =
  \eta_{ab}.
  \tag{12.5.9}\label{eq:12.5.9}
\end{equation}

Now that we have shown how to make any tensor field into a set of coordinate
scalars, we can forget the original tensors \(A^\mu\), \(T_{\mu\nu}\), and so
on, with which we started, and consider how we would construct an action if
we worked from the beginning with the scalars \(A^a\), \(T_{ab}\), and so on.
In this way, a spinor field, like Dirac's electron field, can be brought into
our formalism in precisely the same way as any other field, and its peculiar
Lorentz transformation properties cause no particular trouble. There are now
two invariance principles which must be met in constructing a suitable matter
action \(I_M\).

\begin{description}
  \item[(A)] The action must be generally covariant, with all fields treated
  as coordinate scalars except for the tetrad itself.

  \item[(B)] The Principle of Equivalence requires that special relativity
  should apply in locally inertial frames, and, in particular, that it should
  make no difference which locally inertial frame we choose at each point.
  Thus, since our coordinate-scalar field components \(A^a\), \(T_{ab}\), and
  so on, are defined with respect to an arbitrarily chosen locally inertial
  coordinate system, the field equations and the action must be invariant
  with respect to a redefinition of these locally inertial coordinate systems
  at each point, or, in other words, with respect to Lorentz transformations
  \(\Lambda^a{}_b(x)\) that can depend on position in space-time:
\end{description}
\[
  A^a(x)\longrightarrow\Lambda^a{}_b(x)A^b(x),
\]
\[
  T_{ab}(x)
  \longrightarrow
  \Lambda_a{}^c(x)\Lambda_b{}^d(x)T_{cd}(x),
  \quad\text{etc.},
\]
where
\begin{equation}
  \eta_{ab}\Lambda^a{}_c(x)\Lambda^b{}_d(x)
  =
  \eta_{cd}.
  \tag{12.5.10}\label{eq:12.5.10}
\end{equation}
The tetrad \eqref{eq:12.5.2} changes according to the same rule as \(A^a\):
\begin{equation}
  e^a{}_\mu(x)
  \longrightarrow
  \Lambda^a{}_b(x)e^b{}_\mu(x),
  \tag{12.5.11}\label{eq:12.5.11}
\end{equation}
and in general an arbitrary field \(\psi_n(x)\) will change according to the
rule
\begin{equation}
  \psi_n(x)
  \longrightarrow
  \sum_m[D(\Lambda(x))]_{nm}\psi_m(x),
  \tag{12.5.12}\label{eq:12.5.12}
\end{equation}
where \(D(\Lambda)\) is a matrix representation of the Lorentz group (or at
least of the infinitesimal Lorentz group), of the sort discussed in
Section~\ref{sec:2.12}.

These two invariance principles lead to a dual classification of physical
quantities. A \emph{coordinate scalar} or \emph{coordinate tensor} transforms
as a scalar or a tensor under changes in the coordinate system. A
\emph{Lorentz scalar} or \emph{Lorentz tensor} or \emph{Lorentz spinor}
transforms according to a rule like \eqref{eq:12.5.12}, with
\(D(\Lambda)\) the identity or a tensor representation or a spinor
representation of the infinitesimal Lorentz group, under changes in the
choice of the locally inertial coordinate frame. For instance, a field such
as \eqref{eq:12.5.4} is a coordinate scalar and a Lorentz vector, the Dirac
field of the electron is a coordinate scalar and a Lorentz spinor, and the
tetrad \(e^a{}_\mu\) is a coordinate vector and a Lorentz vector. To be
physically acceptable, the matter action \(I_M\) must be both a coordinate
scalar and a Lorentz scalar.

At this point, the reader may be becoming uneasy. How is the gravitational
field going to get into this sort of theory, when the coordinate-scalar
components \eqref{eq:12.5.9} of the metric tensor are just the constants
\(\eta_{ab}\)? The answer is that gravitational tensor fields appear in the
action because, and only because, of the necessity of introducing
\emph{derivatives} into the theory. If it made sense to construct an action
\(I_M\) solely from fields, and not their derivatives, then it would only be
necessary to choose some arbitrary Lorentz-invariant function
\(\mathcal L(\psi(x))\) of various fields \(\psi_n(x)\) (but not the tetrad),
call them all coordinate scalars, and take the action as
\[
  I_M
  =
  \int\dd^4x\,\sqrt{-g(x)}\,\mathcal L(\psi(x)).
\]
This would then automatically be a coordinate scalar and a Lorentz scalar.
However, the examples discussed in the previous sections of this chapter
show that any physically sensible action must involve derivatives of physical
quantities as well as the quantities themselves. The tetrad field must enter
into the action in such a way as to keep it a coordinate scalar and a Lorentz
scalar despite the presence of derivatives.

An ordinary derivative is of course a coordinate vector, in the sense that
under a coordinate transformation \(x\longrightarrow x'\), it transforms
according to the rule
\[
  \partial_\mu
  \longrightarrow
  \partial'_\mu
  =
  \frac{\partial x^\nu}{\partial x'^\mu}\partial_\nu.
\]
If all the fields appearing in the action were coordinate scalars, there
would be no contravariant indices with which to contract the covariant index
\(\mu\); hence, in order to make the action a coordinate scalar, it is
necessary to introduce the tetrad field, and incorporate derivatives into the
action in the form
\begin{equation}
  e_a{}^\mu\partial_\mu.
  \tag{12.5.13}\label{eq:12.5.13}
\end{equation}
However, although this is a coordinate scalar, it does not have simple
transformation properties under position-dependent Lorentz transformations.
Acting on a general field \(\psi\) with the Lorentz transformation rule
\eqref{eq:12.5.12}, the coordinate-scalar derivative has the transformation
rule
\begin{equation}
\begin{aligned}
  e_a{}^\mu(x)\partial_\mu\psi(x)
  \longrightarrow{}&
  \Lambda_a{}^b(x)e_b{}^\mu(x)
  \partial_\mu\!\left[D(\Lambda(x))\psi(x)\right]
  \\
  ={}&
  \Lambda_a{}^b(x)e_b{}^\mu(x)
  \left\{
    D(\Lambda(x))\partial_\mu\psi(x)
    +[\partial_\mu D(\Lambda(x))]\psi(x)
  \right\}.
\end{aligned}
  \tag{12.5.14}\label{eq:12.5.14}
\end{equation}
However, what we need is to incorporate derivatives into the action in the
form of an operator \(\mathcal D_a\) that is not only a coordinate scalar but
also, unlike \eqref{eq:12.5.13}, is a Lorentz vector, in the sense that for a
position-dependent Lorentz transformation \(\Lambda^a{}_b(x)\),
\begin{equation}
  \mathcal D_a\psi(x)
  \longrightarrow
  \Lambda_a{}^b(x)D(\Lambda(x))\mathcal D_b\psi(x).
  \tag{12.5.15}\label{eq:12.5.15}
\end{equation}
Any action which depends only on various fields \(\psi\) and on their
``derivatives'' \(\mathcal D_a\psi\) will then automatically be independent
of the choice of locally inertial frames if it is invariant under ordinary
constant Lorentz transformations. Inspection of Eq.~\eqref{eq:12.5.14}
shows that we can construct a coordinate-scalar, Lorentz-vector
derivative\hyperref[ch12-ref-1]{\textsuperscript{1}} \(\mathcal D_a\) of the
form
\begin{equation}
  \mathcal D_a
  \equiv
  e_a{}^\mu\left(\partial_\mu+\Gamma_\mu\right),
  \tag{12.5.16}\label{eq:12.5.16}
\end{equation}
where \(\Gamma_\mu\) is a matrix with the Lorentz transformation rule
\begin{equation}
\begin{aligned}
  \Gamma_\mu(x)\longrightarrow{}&
  D(\Lambda(x))\Gamma_\mu(x)D^{-1}(\Lambda(x))
  \\
  &-[\partial_\mu D(\Lambda(x))]D^{-1}(\Lambda(x)).
\end{aligned}
  \tag{12.5.17}\label{eq:12.5.17}
\end{equation}
The inhomogeneous term in Eq.~\eqref{eq:12.5.17} will then cancel the second
term in Eq.~\eqref{eq:12.5.14}, giving \(\mathcal D_a\) the desired
transformation property \eqref{eq:12.5.15}.

In order to determine the structure of the matrix \(\Gamma_\mu(x)\), it will
be sufficient to consider Lorentz transformations that are infinitesimally
close to the identity. Such transformations must be of the form
\begin{equation}
  \Lambda^a{}_b(x)
  =
  \delta^a{}_b+\omega^a{}_b(x),
  \tag{12.5.18}\label{eq:12.5.18}
\end{equation}
with
\begin{equation}
  \omega_{ab}(x)=-\omega_{ba}(x).
  \tag{12.5.19}\label{eq:12.5.19}
\end{equation}
In this case, the matrix \(D\) has the form
\begin{equation}
  D(1+\omega(x))
  =
  1+\frac12\omega^{ab}(x)\sigma_{ab},
  \tag{12.5.20}\label{eq:12.5.20}
\end{equation}
where \(\sigma_{ab}\) are a set of constant matrices that are antisymmetric in
\(a\) and \(b\),
\begin{equation}
  \sigma_{ab}=-\sigma_{ba},
  \tag{12.5.21}\label{eq:12.5.21}
\end{equation}
and that satisfy the commutation relations
\begin{equation}
  [\sigma_{ab},\sigma_{cd}]
  =
  \eta_{cb}\sigma_{ad}
  -\eta_{ca}\sigma_{bd}
  +\eta_{\dd b}\sigma_{ca}
  -\eta_{\dd a}\sigma_{cb}.
  \tag{12.5.22}\label{eq:12.5.22}
\end{equation}
The condition \eqref{eq:12.5.17} now fixes the transformation of
\(\Gamma_\mu(x)\). Under the infinitesimal Lorentz transformation
\eqref{eq:12.5.18}, it transforms according to
\begin{equation}
\begin{aligned}
  \Gamma_\mu(x)\longrightarrow{}&
  \Gamma_\mu(x)
  +\frac12\omega^{ab}(x)[\sigma_{ab},\Gamma_\mu(x)]
  \\
  &-\frac12\sigma_{ab}\partial_\mu\omega^{ab}(x).
\end{aligned}
  \tag{12.5.23}\label{eq:12.5.23}
\end{equation}
Note that \(e^a{}_\nu(x)\) transforms into
\[
  e^a{}_\nu(x)
  \longrightarrow
  e^a{}_\nu(x)+\omega^a{}_b(x)e^b{}_\nu(x),
\]
and therefore, using Eq.~\eqref{eq:12.5.8},
\[
\begin{aligned}
  e_b{}^\nu\nabla_\mu e_{a\nu}
  \longrightarrow{}&
  e_b{}^\nu\nabla_\mu e_{a\nu}
  +\omega_b{}^c e_c{}^\nu\nabla_\mu e_{a\nu}
  \\
  &+\omega_a{}^c e_b{}^\nu\nabla_\mu e_{c\nu}
  +\partial_\mu\omega_{ab}.
\end{aligned}
\]
Multiplying this with \(\sigma^{ab}\) and using the commutation rules
\eqref{eq:12.5.22}, we find that the transformation condition
\eqref{eq:12.5.23} is satisfied by the matrix
\begin{equation}
  \Gamma_\mu(x)
  =
  \frac12\sigma^{ab}e_a{}^\nu(x)\nabla_\mu e_{b\nu}(x).
  \tag{12.5.24}\label{eq:12.5.24}
\end{equation}

To summarize: The effects of gravitation on any physical system can be taken
into account by writing down the matter action or the field equations that
hold in special relativity, and then replacing all derivatives by the
``covariant'' derivatives
\begin{equation}
  \mathcal D_a
  \equiv
  e_a{}^\mu\partial_\mu
  +\frac12\sigma^{bc}e_a{}^\mu e_b{}^\nu
   \nabla_\mu e_{c\nu}.
  \tag{12.5.25}\label{eq:12.5.25}
\end{equation}
This prescription yields a matter action or field equations that are
invariant under general coordinate transformations, with \(e_a{}^\mu\)
regarded as a four-vector and with all other fields regarded as coordinate
scalars, and that also do not depend on how we choose locally inertial frames
in defining the tetrad.

How are we to define the energy-momentum tensor in this formalism? A
variation \(\delta e^a{}_\mu\) in the tetrad field will produce in the matter
action a change
\begin{equation}
  \delta I_M
  =
  \int\dd^4x\,\sqrt{-g}\,
  U_a{}^\mu\,\delta e^a{}_\mu,
  \tag{12.5.26}\label{eq:12.5.26}
\end{equation}
where \(U_a{}^\mu\) is a coordinate vector and a Lorentz vector. Let us
tentatively define the energy-momentum tensor by
\begin{equation}
  T_{\mu\nu}
  \equiv
  e^a{}_\mu U_{a\nu}.
  \tag{12.5.27}\label{eq:12.5.27}
\end{equation}
As required, this is manifestly a coordinate tensor and a Lorentz scalar. To
verify that Eq.~\eqref{eq:12.5.27} is a suitable energy-momentum tensor, we
must also check that it is symmetric,
\begin{equation}
  T_{\mu\nu}=T_{\nu\mu},
  \tag{12.5.28}\label{eq:12.5.28}
\end{equation}
and that it is conserved,
\begin{equation}
  \nabla_\nu T^\nu{}_\lambda=0.
  \tag{12.5.29}\label{eq:12.5.29}
\end{equation}

The symmetry property of the energy-momentum tensor is not at all obvious in
the tetrad formalism, but must be derived from the invariance of the matter
action under the infinitesimal Lorentz transformations
\[
  \Lambda^a{}_b(x)
  =
  \delta^a{}_b+\omega^a{}_b(x),
  \qquad
  \lvert\omega^a{}_b(x)\rvert\ll1.
\]
These transformations will produce small changes in all the dynamic
variables, but the matter action \(I_M\) is supposed to be stationary with
respect to variations in each of these variables except the tetrad, which
enters in \(I_M\) as an external field. Hence we only need to take account of
the change \eqref{eq:12.5.11} in the tetrad field:
\begin{equation}
  \delta e^a{}_\mu(x)
  =
  \omega^a{}_b(x)e^b{}_\mu(x).
  \tag{12.5.30}\label{eq:12.5.30}
\end{equation}
Using Eq.~\eqref{eq:12.5.30} in Eq.~\eqref{eq:12.5.26}, we find that the
invariance of the matter action under position-dependent Lorentz
transformations requires that
\[
  0
  =
  \int\dd^4x\,\sqrt{-g}\,
  U_a{}^\mu e^b{}_\mu\omega^a{}_b.
\]
But \(\omega_{ab}(x)\) is arbitrary, except for the antisymmetry condition
\eqref{eq:12.5.19}, so the coefficient of \(\omega_{ab}(x)\) must be
symmetric:
\[
  U_a{}^\mu e_b{}_\mu
  =
  U_b{}^\mu e_a{}_\mu.
\]
Multiplying this equation by suitable inverse tetrads and using
Eq.~\eqref{eq:12.5.7}, we find that
\begin{equation}
  e^a{}_\lambda U_{a\nu}
  =
  e^a{}_\nu U_{a\lambda},
  \tag{12.5.31}\label{eq:12.5.31}
\end{equation}
which is the same as the expected symmetry condition
\eqref{eq:12.5.28}.

To show that the energy-momentum tensor defined by
Eq.~\eqref{eq:12.5.27} is conserved, we must use the invariance of the matter
action under infinitesimal coordinate transformations:
\[
  x'^\mu=x^\mu+\xi^\mu(x),
  \qquad
  \lvert\xi^\mu\rvert\ \text{very small}.
\]
Such transformations alter the tetrad field by an infinitesimal amount
\begin{equation}
\begin{aligned}
  \delta e^a{}_\mu(x)
  &\equiv
  e'^a{}_\mu(x)-e^a{}_\mu(x)
  \\
  &=
  -e^a{}_\nu(x)\partial_\mu\xi^\nu(x)
  -\partial_\lambda e^a{}_\mu(x)\,\xi^\lambda(x)
  \\
  &=
  -(\Lie_\xi e^a)_\mu.
\end{aligned}
  \tag{12.5.32}\label{eq:12.5.32}
\end{equation}
[Compare Eq.~\eqref{eq:12.3.1}.] Also, all other coordinate-scalar fields
\(\psi(x)\) change by the amounts
\[
  \delta\psi(x)
  \equiv
  \psi'(x)-\psi(x)
  =
  -\partial_\lambda\psi(x)\,\xi^\lambda(x),
\]
but again, the matter action \(I_M\) is stationary with respect to variations
in these fields, so it is only the variation in the tetrad field that matters
here. Using Eq.~\eqref{eq:12.5.32} in Eq.~\eqref{eq:12.5.26}, we find that
the invariance of the matter action \(I_M\) under general coordinate
transformations requires that
\[
  0
  =
  -\int\dd^4x\,\sqrt{-g}\,U_a{}^\mu
  \left[
    e^a{}_\nu\partial_\mu\xi^\nu
    +(\partial_\nu e^a{}_\mu)\xi^\nu
  \right].
\]
But \(\xi^\nu(x)\) is arbitrary, so after integrating by parts we can set the
coefficient of \(\xi^\nu(x)\) equal to zero, and find
\[
  0
  =
  \partial_\mu\!\left(\sqrt{-g}\,U_a{}^\mu e^a{}_\nu\right)
  -\sqrt{-g}\,U_a{}^\mu\partial_\nu e^a{}_\mu.
\]
Using Eqs.~\eqref{eq:12.5.27} and \eqref{eq:12.5.8}, let us write this as
\begin{equation}
  0
  =
  \partial_\mu\!\left(\sqrt{-g}\,T^\mu{}_\nu\right)
  -\sqrt{-g}\,
   T_\rho{}^\mu e_a{}^\rho\partial_\nu e^a{}_\mu.
  \tag{12.5.33}\label{eq:12.5.33}
\end{equation}
According to Eq.~\eqref{eq:12.5.1}, the metric tensor is related to the
tetrad by
\[
  g_{\rho\mu}
  =
  e^a{}_\rho e_{a\mu},
\]
and therefore
\[
  \partial_\nu g_{\rho\mu}
  =
  e_{a\rho}\partial_\nu e^a{}_\mu
  +e_{a\mu}\partial_\nu e^a{}_\rho.
\]
Since \(T_{\rho\mu}\) is symmetric, Eq.~\eqref{eq:12.5.33} may now be written
\begin{equation}
  0
  =
  \partial_\mu\!\left(\sqrt{-g}\,T^\mu{}_\nu\right)
  -\frac12\sqrt{-g}\,T^{\rho\mu}\partial_\nu g_{\rho\mu}.
  \tag{12.5.34}\label{eq:12.5.34}
\end{equation}
But Eqs.~\eqref{eq:3.3.1} and \eqref{eq:4.7.6} give
\[
  \partial_\nu g_{\rho\mu}
  =
  g_{\lambda\mu}\Gamma^\lambda{}_{\nu\rho}
  +g_{\rho\lambda}\Gamma^\lambda{}_{\nu\mu},
  \qquad
  \partial_\mu\ln\sqrt{-g}
  =
  \Gamma^\lambda{}_{\lambda\mu},
\]
so Eq.~\eqref{eq:12.5.34} is the same as the usual conservation law
\eqref{eq:12.5.29}.

Our definition \eqref{eq:12.5.27} of the energy-momentum tensor is thus
completely satisfactory. Note, however, that if the matter action were not
invariant under position-dependent Lorentz transformations, then not only
would \(T_{\mu\nu}\) not be symmetric, but in consequence it also would not
be conserved.

The total action for matter and gravitation is again of the form
\[
  I=I_M+I_G,
\]
with \(I_G\) of the form \eqref{eq:12.4.2}. A variation in the tetrad will
produce in the metric a change given by Eq.~\eqref{eq:12.5.1} as
\[
  \delta g_{\mu\nu}
  =
  e_{a\mu}\delta e^a{}_\nu
  +e_{a\nu}\delta e^a{}_\mu,
\]
so Eq.~\eqref{eq:12.4.3} gives the change in the gravitational action as
\begin{equation}
  \delta I_G
  =
  -\frac1{8\pi G}\int\dd^4x\,\sqrt{-g}\,
  \left(R^\mu{}_\nu-\frac12\delta^\mu{}_\nu R\right)
  e_{a\mu}\delta e^a{}_\nu.
  \tag{12.5.35}\label{eq:12.5.35}
\end{equation}
The total action must be stationary with respect to variations in the tetrad,
so Eqs.~\eqref{eq:12.5.26} and \eqref{eq:12.5.35} yield the field equation
\[
  \left(R^\mu{}_\nu-\frac12\delta^\mu{}_\nu R\right)e_{a\mu}
  =
  8\pi G\,U_{a\nu}.
\]
Contracting this with \(e^a{}_\lambda\), and using
Eqs.~\eqref{eq:12.5.1} and \eqref{eq:12.5.27}, then yields the familiar
Einstein field equation
\begin{equation}
  R_{\lambda\nu}-\frac12g_{\lambda\nu}R
  =
  8\pi G\,T_{\lambda\nu}.
  \tag{12.5.36}\label{eq:12.5.36}
\end{equation}
These equations serve to determine only \(g_{\mu\nu}\), leaving the tetrad
determined only up to a Lorentz transformation \eqref{eq:12.5.11}. However,
the invariance of the matter action under such position-dependent Lorentz
transformations ensures that all tetrads associated with a given metric have
the same physical effects.
